% problem-type: ham-green-mien-chinh-quy
% order: 100
% Nguon: De thi MAT2306 (PTDHR 1), HK2 2024-2025, K67, Cau 1. (co dap an chinh thuc)
% Bo cuc: De vi du -> Cong thuc -> De + Bai giai

\section{Hàm Green \& miền chính quy}

% ---------------------------------------------------------------
\subsection*{Đề bài ví dụ}
% ---------------------------------------------------------------
Xét trụ tứ giác $Q=\{(x,y,z)\in\mathbb{R}^3: z+1>2|y|>2z,\ 0<x<1\}$.
(a) Vẽ $Q$. (b) Chứng minh $Q$ là miền chính quy, từ đó suy ra tồn tại \textbf{duy nhất} một hàm Green đối với $Q$.

\emph{\small(Nguồn: Đề thi MAT2306 --- PTĐHR 1, HK2 2024--2025, K67, Câu 1.)}

% ---------------------------------------------------------------
\subsection*{Nhận ra dạng này khi nào?}
% ---------------------------------------------------------------
\begin{itemize}
  \item Đề hỏi về \textbf{hàm Green}, tính \textbf{chính quy} của miền, hoặc giải được bài Dirichlet.
  \item Thường kèm việc xét từng điểm biên (góc, cạnh) xem có hàm cản không.
\end{itemize}

% ---------------------------------------------------------------
\subsection*{Công thức \& phương pháp}
% ---------------------------------------------------------------
Hai khái niệm cốt lõi ở Phụ lục A: \emph{hàm Green} và \emph{miền chính quy \& hàm cản}. Quy trình:
\begin{enumerate}
  \item Với điểm biên ``trơn'': dùng điều kiện \textbf{hình cầu ngoài} $\Rightarrow$ chính quy.
  \item Với điểm biên ``xấu'' (góc, cạnh): dựng \textbf{hàm cản} cụ thể (thường qua tọa độ trụ/cực).
  \item Mọi điểm biên chính quy $\Rightarrow$ bài Dirichlet $\Delta\Phi=0,\ \Phi=E(X-\cdot)$ trên $\partial Q$ giải được; NLCĐ cho duy nhất $\Rightarrow$ hàm Green tồn tại \& duy nhất.
\end{enumerate}

% ---------------------------------------------------------------
\subsection*{Đề bài \& bài giải}
% ---------------------------------------------------------------
\paragraph{Bước 1 --- điểm biên trơn.} Mọi điểm trên $\partial Q$ \emph{trừ} cạnh
$I=\{(x,0,0):x\in(0,1)\}$ đều thỏa điều kiện \textbf{hình cầu ngoài} (đặt được quả cầu tiếp xúc ngoài)
nên chính quy.

\paragraph{Bước 2 --- cạnh xấu $I$, dựng hàm cản.} Đổi sang tọa độ trụ $x=x,\ y=r\sin\theta,\ z=r\cos\theta$:
\[
  Q=\Big\{0<x<1,\ \tfrac{\pi}{4}<\theta<\tfrac{7\pi}{4},\ 0<r<\tfrac{1}{2|\sin\theta|-\cos\theta}\Big\},\quad
  \Delta u=v_{xx}+v_{rr}+\tfrac1r v_r+\tfrac{1}{r^2}v_{\theta\theta}.
\]
Cạnh $I$ ứng với $r=0$. Hàm
\[
  v=r^{1/2}\sin\!\Big(\frac{\theta-\pi/5}{2}\Big)
\]
là \textbf{hàm cản} tại các điểm của $I$: điều hòa; với $\theta\in(\tfrac{\pi}{4},\tfrac{7\pi}{4})$ thì
$\tfrac{\theta-\pi/5}{2}\in(0,\pi)$ nên $\sin>0$ (dương trong $Q$); và $v\to0$ khi $r\to0$ (triệt tiêu trên $I$).
Vậy mọi điểm biên đều chính quy $\Rightarrow Q$ chính quy.

\paragraph{Bước 3 --- kết luận.}\leavevmode

\begin{nhobox}
Mọi điểm biên chính quy $\Rightarrow$ với mỗi $X\in Q$, bài Dirichlet $\Delta\Phi=0$ trong $Q$,
$\Phi=E(X-\cdot)$ trên $\partial Q$ có nghiệm; NLCĐ $\Rightarrow$ duy nhất. Vậy hàm Green
$G(X,\cdot)=E(X-\cdot)-\Phi$ \textbf{tồn tại và duy nhất}.
\end{nhobox}
